\section*{半双线性型：}

设 A 是一个环， \(\xi \to \overline{\xi}\) 是 A 的一个对合反自同构，即 A 到其自身的一个双射，使得 \(\overline{(\xi + \eta)} = \overline{\xi} + \overline{\eta}\) 、 \(\overline{(\xi\eta)} = \overline{\eta} . \overline{\xi}\) 、 \(\overline{\overline{\xi}} = \xi\) 对 A 中一切 \(\xi, \eta\) 成立。左 A-模 E 上的半双线性型 \(\Phi\) 是指 E × E 到 A 中的一个映射，使得

\[
\begin{array}{c c} \Phi (x + x ^ {\prime}, y) = \Phi (x, y) + \Phi (x ^ {\prime}, y), & \Phi (x, y + y ^ {\prime}) = \Phi (x, y) + \Phi (x, y ^ {\prime}) \\ \Phi (\alpha x, y) = \alpha \Phi (x, y), & \Phi (x, \alpha y) = \Phi (x, y) \bar {\alpha} \end{array}
\]

对 E 中一切 \(\alpha\in\mathbf{A},x,x^{\prime},y,y^{\prime}\) 成立。

设 \(\varepsilon\) 为 A 的中心的一个元素。称 E 上的半双线性型 \(\Phi\) 为 \(\varepsilon\)-埃尔米特型，如果 \(\Phi(y, x) = \varepsilon \overline{\Phi(x, y)}\) 对一切 \(x \in E\) 、 \(y \in E\) 成立；若 \(\varepsilon = 1\) （相应地 \(\varepsilon = -1\) ），则称 \(\Phi\) 为埃尔米特型（相应地，反埃尔米特型）。

当反自同构 \(\xi\to\overline{\xi}\) 为恒等映射（这蕴含环 A 是交换环）时，相应的半双线性型就是双线性型。这时称“对称”以代替“埃尔米特”，称“反对称”以代替“反埃尔米特”（参见第 III 章）。双线性型 \(\Phi\) 若满足 \(\Phi(x,x)=0\) ，则称为交错型；此时它是反对称的，而当 A 是特征 \(\neq2\) 的域时，其逆亦真（参见第 III 章）。

称 \(\varepsilon\)-埃尔米特型满足条件 (T)，如果对一切 \(x \in E\) ，存在 \(\lambda \in A\) 使 \(\Phi(x, x) = \lambda + \varepsilon \overline{\lambda}\) 。当 \(\varepsilon = 1\) 且 \(A\) 是特征 \(\neq 2\) 的域，或当 \(\Phi\) 是交错型时，这一条件总得到满足。

